% Options for packages loaded elsewhere
\PassOptionsToPackage{unicode}{hyperref}
\PassOptionsToPackage{hyphens}{url}
\documentclass[
  ignorenonframetext,
]{beamer}
\newif\ifbibliography
\usepackage{pgfpages}
\setbeamertemplate{caption}[numbered]
\setbeamertemplate{caption label separator}{: }
\setbeamercolor{caption name}{fg=normal text.fg}
\beamertemplatenavigationsymbolsempty
% remove section numbering
\setbeamertemplate{part page}{
  \centering
  \begin{beamercolorbox}[sep=16pt,center]{part title}
    \usebeamerfont{part title}\insertpart\par
  \end{beamercolorbox}
}
\setbeamertemplate{section page}{
  \centering
  \begin{beamercolorbox}[sep=12pt,center]{section title}
    \usebeamerfont{section title}\insertsection\par
  \end{beamercolorbox}
}
\setbeamertemplate{subsection page}{
  \centering
  \begin{beamercolorbox}[sep=8pt,center]{subsection title}
    \usebeamerfont{subsection title}\insertsubsection\par
  \end{beamercolorbox}
}
% Prevent slide breaks in the middle of a paragraph
\widowpenalties 1 10000
\raggedbottom
\AtBeginPart{
  \frame{\partpage}
}
\AtBeginSection{
  \ifbibliography
  \else
    \frame{\sectionpage}
  \fi
}
\AtBeginSubsection{
  \frame{\subsectionpage}
}
\usepackage{iftex}
\ifPDFTeX
  \usepackage[T1]{fontenc}
  \usepackage[utf8]{inputenc}
  \usepackage{textcomp} % provide euro and other symbols
\else % if luatex or xetex
  \usepackage{unicode-math} % this also loads fontspec
  \defaultfontfeatures{Scale=MatchLowercase}
  \defaultfontfeatures[\rmfamily]{Ligatures=TeX,Scale=1}
\fi
\usepackage{lmodern}
\ifPDFTeX\else
  % xetex/luatex font selection
\fi
% Use upquote if available, for straight quotes in verbatim environments
\IfFileExists{upquote.sty}{\usepackage{upquote}}{}
\IfFileExists{microtype.sty}{% use microtype if available
  \usepackage[]{microtype}
  \UseMicrotypeSet[protrusion]{basicmath} % disable protrusion for tt fonts
}{}
\makeatletter
\@ifundefined{KOMAClassName}{% if non-KOMA class
  \IfFileExists{parskip.sty}{%
    \usepackage{parskip}
  }{% else
    \setlength{\parindent}{0pt}
    \setlength{\parskip}{6pt plus 2pt minus 1pt}}
}{% if KOMA class
  \KOMAoptions{parskip=half}}
\makeatother
\usepackage{longtable,booktabs,array}
\newcounter{none} % for unnumbered tables
\usepackage{calc} % for calculating minipage widths
\usepackage{caption}
% Make caption package work with longtable
\makeatletter
\def\fnum@table{\tablename~\thetable}
\makeatother
\setlength{\emergencystretch}{3em} % prevent overfull lines
\providecommand{\tightlist}{%
  \setlength{\itemsep}{0pt}\setlength{\parskip}{0pt}}
\usepackage{bookmark}
\IfFileExists{xurl.sty}{\usepackage{xurl}}{} % add URL line breaks if available
\urlstyle{same}
\hypersetup{
  hidelinks,
  pdfcreator={LaTeX via pandoc}}

\author{\texorpdfstring{}{}}
\date{}

\begin{document}

\section{Performance targets and measured runs on packaged
fixtures}\label{performance-targets-and-measured-runs-on-packaged-fixtures}

\begin{frame}[fragile]{Performance characteristics}
\protect\phantomsection\label{performance-characteristics}
The architecture targets the following benchmarks on a 4-core machine,
as specified in \texttt{../cogant/docs/architecture/README.md}:

{\def\LTcaptype{none} % do not increment counter
\begin{longtable}[]{@{}lll@{}}
\toprule\noalign{}
Repository size & Target wall-clock time & Memory budget \\
\midrule\noalign{}
\endhead
10K functions & \textless{} 30 s & \textless{} 500 MB \\
100K functions & \textless{} 5 min & \textless{} 2 GB \\
1M functions & \textless{} 1 hr & \textless{} 2 GB (streaming) \\
\bottomrule\noalign{}
\end{longtable}
}

These are architecture targets, not benchmark claims from this
manuscript. They assume the Python orchestration layer with Rust
acceleration on critical paths (graph construction, rule matching, and
Generalized Notation Notation section/tensor packing in
\texttt{cogant-gnn}). In the current v0.5.x release, Rust acceleration
is partially wired --- \texttt{cogant.\_rust} exposes a PyO3
\texttt{connected\_components} FFI for graph construction behind the
\texttt{COGANT\_USE\_RUST} feature flag --- and a pure-Python fallback
handles the remaining code paths.

Current \texttt{PipelineRunner} behavior is stage-sequential with
per-stage error capture and continuation. It does not currently expose
built-in incremental checkpoint/resume in \texttt{cogant.api.pipeline};
treat checkpointing as a potential outer-orchestration feature rather
than a guaranteed package-level runtime behavior.
\end{frame}

\begin{frame}[fragile]{Measured runs on packaged fixtures}
\protect\phantomsection\label{measured-runs-on-packaged-fixtures}
The following tables record measurements taken by running the shipped
\texttt{RoundtripOrchestrator}
(\texttt{../cogant/examples/orchestrate\_roundtrip.py}) against every
fixture distributed with the package. Three fixtures are the
control-positive synthetic repositories under
\texttt{../cogant/examples/control\_positive/} (\texttt{calculator},
\texttt{event\_pipeline}, \texttt{flask\_mini}); the other three are
real-world code under \texttt{../cogant/examples/real\_world/}
(\texttt{flask\_app}, a six-module Flask service;
\texttt{requests\_lib}, a six-module reduction of the \texttt{requests}
HTTP library; and \texttt{json\_stdlib}, a four-module reduction of the
CPython \texttt{json} package). Each run executes the full static
pipeline (ingest, parse, symbols, imports, call graph, program graph,
translation, state-space compilation, GNN package build, validation).
Wall-clock times were measured on a single macOS workstation with the
pure-Python fallback implementations --- the v0.5.x PyO3
\texttt{connected\_components} FFI is disabled for this canonical run
(\texttt{COGANT\_USE\_RUST=0}) so that these numbers correspond to the
Python orchestration layer with no native crates loaded.

All numbers in Tables 4--7 are regenerated by
\texttt{../cogant/evaluation/figures/generate\_figures.py}, which
re-runs the orchestrator over every fixture, re-reads the emitted JSON,
and writes \texttt{../cogant/evaluation/figures/metrics.json} alongside
the figure PNGs. Structural metrics (nodes, edges, edge-kind and
node-kind breakdowns, LOC, file counts) are deterministic; rule-driven
metrics (total mappings, state variables, observations, actions,
transitions) vary by at most one or two units across runs because the
extractor walks dictionaries whose ordering is process-local, and
wall-clock times vary by a few seconds depending on whether the
visualization pass rasterizes PNGs. The figures below match the
canonical \texttt{metrics.json} committed under
\texttt{../cogant/evaluation/figures/}.

\textbf{Table 4. Repository-level pipeline metrics (canonical run,
COGANT v0.1.0).}

{\def\LTcaptype{none} % do not increment counter
\begin{longtable}[]{@{}
  >{\raggedright\arraybackslash}p{(\linewidth - 24\tabcolsep) * \real{0.0588}}
  >{\raggedleft\arraybackslash}p{(\linewidth - 24\tabcolsep) * \real{0.0784}}
  >{\raggedleft\arraybackslash}p{(\linewidth - 24\tabcolsep) * \real{0.0784}}
  >{\raggedleft\arraybackslash}p{(\linewidth - 24\tabcolsep) * \real{0.0784}}
  >{\raggedleft\arraybackslash}p{(\linewidth - 24\tabcolsep) * \real{0.0784}}
  >{\raggedleft\arraybackslash}p{(\linewidth - 24\tabcolsep) * \real{0.0784}}
  >{\raggedleft\arraybackslash}p{(\linewidth - 24\tabcolsep) * \real{0.0784}}
  >{\raggedleft\arraybackslash}p{(\linewidth - 24\tabcolsep) * \real{0.0784}}
  >{\raggedleft\arraybackslash}p{(\linewidth - 24\tabcolsep) * \real{0.0784}}
  >{\raggedleft\arraybackslash}p{(\linewidth - 24\tabcolsep) * \real{0.0784}}
  >{\raggedleft\arraybackslash}p{(\linewidth - 24\tabcolsep) * \real{0.0784}}
  >{\raggedleft\arraybackslash}p{(\linewidth - 24\tabcolsep) * \real{0.0784}}
  >{\raggedleft\arraybackslash}p{(\linewidth - 24\tabcolsep) * \real{0.0784}}@{}}
\toprule\noalign{}
\begin{minipage}[b]{\linewidth}\raggedright
Fixture
\end{minipage} & \begin{minipage}[b]{\linewidth}\raggedleft
Files
\end{minipage} & \begin{minipage}[b]{\linewidth}\raggedleft
LOC
\end{minipage} & \begin{minipage}[b]{\linewidth}\raggedleft
Nodes
\end{minipage} & \begin{minipage}[b]{\linewidth}\raggedleft
Edges
\end{minipage} & \begin{minipage}[b]{\linewidth}\raggedleft
Mappings
\end{minipage} & \begin{minipage}[b]{\linewidth}\raggedleft
State vars
\end{minipage} & \begin{minipage}[b]{\linewidth}\raggedleft
Obs
\end{minipage} & \begin{minipage}[b]{\linewidth}\raggedleft
Actions
\end{minipage} & \begin{minipage}[b]{\linewidth}\raggedleft
Transitions
\end{minipage} & \begin{minipage}[b]{\linewidth}\raggedleft
GNN sections
\end{minipage} & \begin{minipage}[b]{\linewidth}\raggedleft
GNN score
\end{minipage} & \begin{minipage}[b]{\linewidth}\raggedleft
Wall-clock (s)
\end{minipage} \\
\midrule\noalign{}
\endhead
\texttt{calculator} & 1 & 122 & 12 & 27 & 5 & 0 & 3 & 1 & 1 & 31 & 100.0
& 7.42 \\
\texttt{event\_pipeline} & 1 & 147 & 23 & 66 & 20 & 1 & 9 & 10 & 10 & 31
& 100.0 & 8.58 \\
\texttt{flask\_mini} & 1 & 168 & 26 & 51 & 19 & 3 & 2 & 14 & 14 & 31 &
100.0 & 8.80 \\
\texttt{flask\_app} & 6 & 853 & 98 & 597 & 51 & 16 & 19 & 16 & 16 & 31 &
100.0 & 14.58 \\
\texttt{requests\_lib} & 6 & 750 & 98 & 345 & 46 & 9 & 31 & 7 & 7 & 31 &
100.0 & 12.19 \\
\texttt{json\_stdlib} & 4 & 1231 & 29 & 68 & 8 & 3 & 5 & 0 & 0 & 31 &
100.0 & 9.08 \\
\bottomrule\noalign{}
\end{longtable}
}

``GNN sections'' counts the level-two Markdown headings emitted in
\texttt{model.gnn.md}, which on every fixture sits at 31 (the 18 core
Generalized Notation Notation sections plus section-specific subheadings
for state space, observations, actions, and transitions). ``GNN score''
is the \texttt{score} field returned by \texttt{GNNValidator.validate()}
on the compiled \texttt{gnn\_package/} directory; every fixture
validates at 100.0 with zero errors and zero warnings. The six fixtures
together cover one-to-six file repositories and 122 to 1231 lines of
code, exercising the pipeline on both minimal control positives and
small real-world modules. Wall-clock times fall between roughly seven
and fifteen seconds on a 2024-class Apple-silicon workstation; the bulk
of the cost on the larger fixtures is the call-graph construction step
in \texttt{CallGraphBuilder} plus the PNG rasterization pass in
\texttt{cogant.viz.png\_export}.

\textbf{Table 5. Program graph composition by fixture.}

Node kinds (MODULE / CLASS / METHOD / FUNCTION) and edge kinds (CONTAINS
/ WRITES / READS / CALLS / IMPORTS / INHERITS) are populated directly
from the AST extractor and call-graph builder; all other kinds listed in
\texttt{cogant.schemas.core.NodeKind} and \texttt{EdgeKind} remain
unused on these fixtures because the Python front end currently focuses
on the structural core. The counts below are taken verbatim from the
\texttt{statistics} block of each fixture's
\texttt{program\_graph.json}.

{\def\LTcaptype{none} % do not increment counter
\begin{longtable}[]{@{}
  >{\raggedright\arraybackslash}p{(\linewidth - 20\tabcolsep) * \real{0.0698}}
  >{\raggedleft\arraybackslash}p{(\linewidth - 20\tabcolsep) * \real{0.0930}}
  >{\raggedleft\arraybackslash}p{(\linewidth - 20\tabcolsep) * \real{0.0930}}
  >{\raggedleft\arraybackslash}p{(\linewidth - 20\tabcolsep) * \real{0.0930}}
  >{\raggedleft\arraybackslash}p{(\linewidth - 20\tabcolsep) * \real{0.0930}}
  >{\raggedleft\arraybackslash}p{(\linewidth - 20\tabcolsep) * \real{0.0930}}
  >{\raggedleft\arraybackslash}p{(\linewidth - 20\tabcolsep) * \real{0.0930}}
  >{\raggedleft\arraybackslash}p{(\linewidth - 20\tabcolsep) * \real{0.0930}}
  >{\raggedleft\arraybackslash}p{(\linewidth - 20\tabcolsep) * \real{0.0930}}
  >{\raggedleft\arraybackslash}p{(\linewidth - 20\tabcolsep) * \real{0.0930}}
  >{\raggedleft\arraybackslash}p{(\linewidth - 20\tabcolsep) * \real{0.0930}}@{}}
\toprule\noalign{}
\begin{minipage}[b]{\linewidth}\raggedright
Fixture
\end{minipage} & \begin{minipage}[b]{\linewidth}\raggedleft
MODULE
\end{minipage} & \begin{minipage}[b]{\linewidth}\raggedleft
CLASS
\end{minipage} & \begin{minipage}[b]{\linewidth}\raggedleft
METHOD
\end{minipage} & \begin{minipage}[b]{\linewidth}\raggedleft
FUNCTION
\end{minipage} & \begin{minipage}[b]{\linewidth}\raggedleft
CONTAINS
\end{minipage} & \begin{minipage}[b]{\linewidth}\raggedleft
WRITES
\end{minipage} & \begin{minipage}[b]{\linewidth}\raggedleft
READS
\end{minipage} & \begin{minipage}[b]{\linewidth}\raggedleft
CALLS
\end{minipage} & \begin{minipage}[b]{\linewidth}\raggedleft
IMPORTS
\end{minipage} & \begin{minipage}[b]{\linewidth}\raggedleft
INHERITS
\end{minipage} \\
\midrule\noalign{}
\endhead
\texttt{calculator} & 1 & 1 & 10 & --- & 11 & 5 & 9 & 2 & --- & --- \\
\texttt{event\_pipeline} & 1 & 6 & 16 & --- & 22 & 1 & 10 & 30 & --- &
3 \\
\texttt{flask\_mini} & 1 & 7 & 18 & --- & 25 & 6 & 7 & 11 & --- & 2 \\
\texttt{flask\_app} & 6 & 25 & 57 & 10 & 92 & 15 & 38 & 433 & 10 & 9 \\
\texttt{requests\_lib} & 6 & 20 & 59 & 13 & 92 & 8 & 42 & 183 & 10 &
10 \\
\texttt{json\_stdlib} & 4 & 3 & 9 & 13 & 25 & 3 & 6 & 32 & 2 & --- \\
\bottomrule\noalign{}
\end{longtable}
}

The distribution matches the intuitive shape of the fixtures:
class-heavy repositories (\texttt{flask\_app}, \texttt{requests\_lib},
\texttt{flask\_mini}) show METHOD-dominated node counts and INHERITS
edges, while the functional \texttt{json\_stdlib} shows a balanced
METHOD/FUNCTION split and no INHERITS edges. CALLS edges dominate the
large repositories (\texttt{flask\_app} at 433, \texttt{requests\_lib}
at 183), confirming that the call-graph construction step in
\texttt{CallGraphBuilder} is the primary source of edge density on real
code.

\textbf{Table 6. State-space compilation outputs.}

For each fixture the \texttt{StateSpaceCompiler} emits a
\texttt{StateSpaceModel} whose variables, observations, actions, and
transitions are then packaged into
\texttt{gnn\_package/state\_space.json}, \texttt{observations.json},
\texttt{actions.json}, and \texttt{transitions.json}. The counts reflect
the end-to-end behavior of the compiler on these inputs, not the rule
engine's raw mapping output.

{\def\LTcaptype{none} % do not increment counter
\begin{longtable}[]{@{}
  >{\raggedright\arraybackslash}p{(\linewidth - 10\tabcolsep) * \real{0.1304}}
  >{\raggedleft\arraybackslash}p{(\linewidth - 10\tabcolsep) * \real{0.1739}}
  >{\raggedleft\arraybackslash}p{(\linewidth - 10\tabcolsep) * \real{0.1739}}
  >{\raggedleft\arraybackslash}p{(\linewidth - 10\tabcolsep) * \real{0.1739}}
  >{\raggedleft\arraybackslash}p{(\linewidth - 10\tabcolsep) * \real{0.1739}}
  >{\raggedleft\arraybackslash}p{(\linewidth - 10\tabcolsep) * \real{0.1739}}@{}}
\toprule\noalign{}
\begin{minipage}[b]{\linewidth}\raggedright
Fixture
\end{minipage} & \begin{minipage}[b]{\linewidth}\raggedleft
State variables
\end{minipage} & \begin{minipage}[b]{\linewidth}\raggedleft
Observations
\end{minipage} & \begin{minipage}[b]{\linewidth}\raggedleft
Actions
\end{minipage} & \begin{minipage}[b]{\linewidth}\raggedleft
Transitions
\end{minipage} & \begin{minipage}[b]{\linewidth}\raggedleft
Policies
\end{minipage} \\
\midrule\noalign{}
\endhead
\texttt{calculator} & 0 & 3 & 1 & 1 & 1 \\
\texttt{event\_pipeline} & 1 & 9 & 10 & 10 & 4 \\
\texttt{flask\_mini} & 3 & 2 & 14 & 14 & 6 \\
\texttt{flask\_app} & 16 & 19 & 16 & 16 & 6 \\
\texttt{requests\_lib} & 9 & 31 & 7 & 7 & 1 \\
\texttt{json\_stdlib} & 3 & 5 & 0 & 0 & 1 \\
\bottomrule\noalign{}
\end{longtable}
}

\texttt{calculator} compiles zero hidden-state variables because the
fixture exposes pure arithmetic methods whose WRITES-edge footprint does
not cross the classifier threshold used by
\texttt{StateVariableExtractor}; it still compiles a single action and a
single transition, so the pipeline remains end-to-end valid and the GNN
package still validates at 100.0. \texttt{json\_stdlib} compiles zero
actions because the reduced CPython \texttt{json} sources consist almost
entirely of function-level utilities whose semantic role falls outside
the ACTION-mapping rule set, but the pipeline still emits a compiled
\texttt{actions.json} with an empty \texttt{actions} list and a single
default policy --- downstream consumers see the same schema as for the
action-bearing fixtures. The \texttt{requests\_lib} fixture has the
highest observation count in the table because its session/adapter
classes expose a large number of read-only attributes that the rule
engine matches as \texttt{OBSERVATION} mappings.

\textbf{Table 7. Output artifacts per run.}

Every fixture emits the same \texttt{gnn\_package/} directory layout
with 19 canonical files: \texttt{model.gnn.md}, \texttt{model.gnn.json},
the section JSONs (\texttt{state\_space}, \texttt{observations},
\texttt{actions}, \texttt{actions\_policies}, \texttt{transitions},
\texttt{preferences}, \texttt{preferences\_constraints},
\texttt{factors}, \texttt{connections}, \texttt{ontology},
\texttt{provenance}), the \texttt{program\_graph.json} and
\texttt{process\_model.json} snapshots, the
\texttt{markov\_blanket.json} and \texttt{markov\_network.json}
extracts, and the \texttt{manifest.json} index. The
\texttt{RoundtripOrchestrator} additionally writes diagnostic artifacts
at the top level of \texttt{output\_dir}: six Mermaid diagrams (class,
state, sequence, dependency, boundary, semantic flow), typed graph and
Cytoscape exports, a GraphML file, a Parquet table, a dashboard HTML
site, simulation traces, and a GNN execution report.

{\def\LTcaptype{none} % do not increment counter
\begin{longtable}[]{@{}
  >{\raggedright\arraybackslash}p{(\linewidth - 6\tabcolsep) * \real{0.2000}}
  >{\raggedleft\arraybackslash}p{(\linewidth - 6\tabcolsep) * \real{0.2667}}
  >{\raggedleft\arraybackslash}p{(\linewidth - 6\tabcolsep) * \real{0.2667}}
  >{\raggedleft\arraybackslash}p{(\linewidth - 6\tabcolsep) * \real{0.2667}}@{}}
\toprule\noalign{}
\begin{minipage}[b]{\linewidth}\raggedright
Fixture
\end{minipage} & \begin{minipage}[b]{\linewidth}\raggedleft
\texttt{gnn\_package/} files
\end{minipage} & \begin{minipage}[b]{\linewidth}\raggedleft
Validation errors
\end{minipage} & \begin{minipage}[b]{\linewidth}\raggedleft
Validation warnings
\end{minipage} \\
\midrule\noalign{}
\endhead
\texttt{calculator} & 19 & 0 & 0 \\
\texttt{event\_pipeline} & 19 & 0 & 0 \\
\texttt{flask\_mini} & 19 & 0 & 0 \\
\texttt{flask\_app} & 19 & 0 & 0 \\
\texttt{requests\_lib} & 19 & 0 & 0 \\
\texttt{json\_stdlib} & 19 & 0 & 0 \\
\bottomrule\noalign{}
\end{longtable}
}

These numbers were collected by the reproducible script
\texttt{../cogant/evaluation/figures/generate\_figures.py}, which
re-runs the orchestrator over every fixture, re-reads the emitted JSON,
and writes both the summary table and the figures used in this
manuscript into \texttt{../cogant/evaluation/figures/} (namely
\texttt{fig1\_graph\_sizes.png}, \texttt{fig2\_node\_kinds.png},
\texttt{fig3\_state\_space.png}, and
\texttt{fig4\_pipeline\_latency.png}, plus the machine-readable
\texttt{metrics.json}).
\end{frame}

\end{document}
